% In this file you should put the actual content of the blueprint.
% It will be used both by the web and the print version.
%
% Blueprint of Orig2 --- the rebuild library of the Klondike
% macro-state program (lean-model/Orig2), drawn from the live
% cards: MAP.md, the umbrella docstrings, and the source heads.
% Every \lean name was verified against the Orig2 sources.

\chapter{The program}
\label{chap:program}

This blueprint maps \texttt{Orig2}, the rebuild library of the
Klondike macro-state program, and the campaign that is running on
top of it.

The game is Klondike solitaire compressed to its macro-state
physics.  A \emph{state} is a full board: foundations, tableau
piles, stock and waste.  A \emph{witness journey} is a play in
which every letter is reversible.  The \emph{macro class} of a
state is the set of states joined to it by witness journeys both
ways.  The program's question is the \emph{saturation ceiling}:
when the board admits no further reversible progress, how much of
the foundation layout is forced by the class alone?

The certified substrate (Chapter~\ref{chap:substrate}) carries
the game model, the fiber transfer, and the decidability bridge.
The fork theory (Chapter~\ref{chap:fork}) develops the local
obstruction that any class-level difference must pass through.
The fiber hierarchy (Chapter~\ref{chap:fiber}) assembles the
dichotomy family and the campaign payoff.  The boundary cases
(Chapter~\ref{chap:boundary}) are the counterexamples that fix the
target shapes.  The campaign target (Chapter~\ref{chap:campaign})
is the open summit \texttt{SatDifferFork}; the drain and the
asymmetry bridge are closed, the remaining obligations are the
bracket existence at the minimality gate and the split block, and
fork-position rigidity is closed unconditionally.

\paragraph{How to read the graph.}  A green border means the
statement is formalized in Lean; a light green fill means the
proof is formalized; a light blue fill means the proof is ready to
be formalized (all prerequisites done); a dark green fill means
the result and all its ancestors are formalized.  The open
summit is exactly the node with a green border and no proof fill
at the top of the graph.

\chapter{The substrate}
\label{chap:substrate}

\section{The game model}

\begin{definition}[The game model]\label{def:model}
\lean{State, State.WF, rungAt, RungEq}
\leanok
A state of the game is the full board.  A state is \emph{well
formed} (\texttt{WF}) when every zone obeys the conservation laws.
The \emph{rung} $\mathrm{rungAt}(st, \sigma)$ of a suit $\sigma$
is the height of its foundation row; two states are
\emph{rung-equal} when all four rungs agree.
\end{definition}

\begin{definition}[The reversible orbit]\label{def:orbit}
\lean{ShufflePlayW, RevEqW, sameOrbitSetoid}
\leanok
\uses{def:model}
A \emph{witness shuffle} is a play whose every letter is
reversible at the weight grade.  Two states are \emph{same-orbit}
mates when witness shuffles join them both ways; this is the
macro class of a state.
\end{definition}

\begin{definition}[S-saturation]\label{def:ssaturated}
\lean{CanRaiseS, SSaturated}
\leanok
\uses{def:model}
A card is \emph{S-raisable} when it can be raised outright or
after any cascade of reversible slides.  A state is
\emph{S-saturated} when no card is S-raisable: the board admits
no further reversible progress.
\end{definition}

\begin{definition}[Move fibers]\label{def:fibers}
\lean{MoveFiber, CargoSpan, FlatSpan, MacroSpan}
\leanok
\uses{def:orbit}
The two-way closure of a move set $M$ is the \emph{move fiber}
$\mathrm{MoveFiber}(M)$.  The cargo moves give the cargo span,
the flat moves the flat span, and all moves the macro span, with
inclusions cargo $\subset$ flat $\subset$ macro.
\end{definition}

\begin{definition}[The position grid]\label{def:position}
\lean{WFState, position_meet}
\leanok
\uses{def:fibers, def:model}
The readout quotient by rung equality is the rung class.  The
\emph{Position} of a well-formed state is its macro span paired
with its rung class: the meet of the move-generated and readout
axes.
\end{definition}

\section{The fiber transfer}

The transfer theorem is the substrate's crown: the original ledger
stub, discharged verbatim.  Its route passes through the pair-free
stepping stone, the first-touch restart, and the non-king/king
split.

\begin{theorem}[The pair-free transfer]\label{thm:pairFree}
\lean{swappableIn_transfer_of_pairFreePath}
\leanok
\uses{def:orbit}
If the class journey never seats the twin pair, the strict swap
transfers along the class: a swappable pair at the source state
stays swappable at every class-mate with the same foundations.
\end{theorem}

\begin{proof}
\leanok
Discharged in the Orig2 sources --- the card's axiom battery prints [propext, Quot.sound].
\end{proof}

\begin{theorem}[First-touch restart]\label{thm:firstTouch}
\lean{firstTouchRestart, FirstTouchRestart}
\leanok
\uses{thm:pairFree}
The first-touch restart property holds: the class journey can be
restarted at the first touch of either twin, so the pair-free
play splices onto the run splices.
\end{theorem}

\begin{proof}
\leanok
Discharged in the Orig2 sources --- the card's axiom battery prints [propext, Quot.sound].
\end{proof}

\begin{theorem}[The last-return route]\label{thm:lastReturn}
\lean{fiber_transfer_of_last_return}
\leanok
\uses{thm:firstTouch}
The corrected fiber transfer via the last return: every class
journey can be re-threaded so that the pair-free reading applies
at the last return of the twin.
\end{theorem}

\begin{proof}
\leanok
Discharged in the Orig2 sources --- the card's axiom battery prints [propext, Quot.sound].
\end{proof}

\begin{theorem}[The non-king transfer]\label{thm:nonking}
\lean{fiber_transfer_of_nonking}
\leanok
\uses{thm:lastReturn}
The transfer for a non-king pair: the non-king twin cannot hide
(its seat pins it), so the strict swap transfers.
\end{theorem}

\begin{proof}
\leanok
Discharged in the Orig2 sources --- the card's axiom battery prints [propext, Quot.sound].
\end{proof}

\begin{theorem}[The king transfer]\label{thm:king}
\lean{fiber_transfer_of_king}
\leanok
\uses{thm:lastReturn}
The transfer for a king pair: a swappable king is exactly a
vacancy witness, and the vacancy transfers along the class.
\end{theorem}

\begin{proof}
\leanok
Discharged in the Orig2 sources --- the card's axiom battery prints [propext, Quot.sound].
\end{proof}

\begin{theorem}[The fiber transfer, discharged]\label{thm:transferFull}
\lean{fiber_transfer_full}
\leanok
\uses{thm:twinMembership}
The original ledger stub's statement, verbatim: well-formed
same-orbit states with equal foundations and a strict witness
swap at the source carry a strict witness swap to every class-mate
that holds both twins.  After the tower collapse this is a
one-line delegate: the seat-free strengthening beneath it proved
strong enough to carry the crown directly, and the rank split
lives in its own skeleton home.
\end{theorem}

\begin{proof}
\leanok
\uses{thm:twinMembership}
Exact the seat-free strengthening --- the tower collapsed when
the seat dichotomy (both-or-neither) made the pile-holding
premises redundant.
\end{proof}

\begin{theorem}[The twin-membership transfer]\label{thm:twinMembership}
\lean{fiber_transfer_twin_membership, fiber_transfer_seat_free, fiber_transfer_founds_necessary}
\leanok
\uses{thm:nonking, thm:king}
The companion strengthenings: the swap witness transfers along
the class whenever the twins' foundation-membership is preserved
--- the first-seat-change argument, since a twin's only exits
are a raise or a slide onto an exposed fitting top, and the
fitting class is exactly the two parents.  The seat premises
are redundant (both-or-neither), and the founds rows are
necessary.  This closes the depth-0 rung of the
ladder-completeness residue.
\end{theorem}

\begin{proof}
\leanok
Discharged in the Orig2 card (green, axiom battery [propext, Quot.sound]).
\end{proof}

\begin{theorem}[The stratified quotient's universal property]\label{thm:stratifiedQuotient}
\lean{stratOrbit, swappable_stratum_invariant, SwappableStratum, swappableStratum_mk, stratOrbit_universal, swappableStratum_unique, swappable_of_pairFreePlay}
\leanok
\uses{thm:twinMembership}
The transfer, generalized: within a fiber of the map taking a
well-formed state to its reversible class and its twins'
membership profile, swappability is constant --- so
\texttt{SwappableIn} descends to the stratified quotient,
uniquely, and so does every stratum-invariant predicate.  The
two house axioms become load-bearing by design:
\texttt{Quot.sound} is the quotient's universal property.  This
discharges the quotient-formalization residue.  The unified
pair-free transfer rides the same stratum constancy: no seat
premises, both directions, subsuming the engine-grade pair-free
shape for every external consumer --- and the supersession
survey's verdict is that the pair-free family is topologically
un-supersessable, the stratum kit sitting above it by DAG
necessity.
\end{theorem}

\begin{proof}
\leanok
Discharged in the Orig2 card (green, axiom battery [propext, Quot.sound]).
\end{proof}

\begin{theorem}[The stratification link]\label{thm:stratificationLink}
\lean{stratOrbit_of_class_rungs, swappable_of_class_rungs, stratOrbit_of_bareFlat, swappable_of_bareFlat}
\leanok
\uses{def:position, def:fibers, thm:stratifiedQuotient}
Positions refine strata, and so do flat fibers: two states with
the same class and the same rung vector --- the kernel of the
Position meet --- lie in one stratum, since the rungs pin the
founds and the founds pin the twins' membership profile; and a
flat path never moves a foundation row, so each flat fiber lies
in a single stratum.  Hence swappability is both
position-determined and flat-fiber-determined: the stratified
quotient sees strictly coarser than both live axes of the
campaign.
\end{theorem}

\begin{proof}
\leanok
Discharged in the Orig2 card (green, axiom battery [propext, Quot.sound]).
\end{proof}

\begin{theorem}[The decidability bridge]\label{thm:decidable}
\lean{canRaiseS_iff, canRaiseSB}
\leanok
\uses{def:ssaturated}
S-raiseability is a flat, decidable predicate: it is equivalent
to the next-up flag plus a finite seat scan, and the scan is
certified by a Bool checker with a proved correctness iff.
\end{theorem}

\begin{proof}
\leanok
Discharged in the Orig2 sources --- the card's axiom battery prints [propext, Quot.sound].
\end{proof}

\begin{theorem}[The free-card invariant]\label{thm:freeorbit}
\lean{freeCardS_orbit_rung, canRaiseS_orbit_rung}
\leanok
\uses{def:ssaturated, def:orbit}
Being a free card is an orbit-and-rung invariant: the macro class
and the rung vector decide it.  Likewise S-raiseability itself is
an orbit-and-rung invariant at well-formed states.
\end{theorem}

\begin{proof}
\leanok
Discharged in the Orig2 sources --- the card's axiom battery prints [propext, Quot.sound].
\end{proof}

\chapter{The fork theory}
\label{chap:fork}

\section{Detection and the critical pair}

\begin{definition}[Fork detection]\label{def:fork}
\lean{forkState, ForkWitness}
\leanok
\uses{def:ssaturated}
A state is a \emph{fork} when some card $x$ has its S-live pair
exactly $\{x, x.\mathrm{twin}\}$ while the twin is blocked in the
vanilla grade: the pair-saturation iff plus the orientation pin.
\texttt{ForkWitness} is the per-witness parcel.
\end{definition}

\begin{theorem}[The two-raise dichotomy]\label{thm:criticalPair}
\lean{twin_raise_square, canRaise_or_covered, forkWitness_iff_covered}
\leanok
\uses{def:fork}
At a fork the witness is exactly the covered card: either the
raise of the one member or the cover of the twin, and the
witness-iff-covered equivalence pins the fork to its cover.
\end{theorem}

\begin{proof}
\leanok
Discharged in the Orig2 sources --- the card's axiom battery prints [propext, Quot.sound].
\end{proof}

\begin{theorem}[The rung obstruction]\label{thm:rungObstruction}
\lean{arms_never_flatJoin, join_of_not_covered, no_raise_after_twin_raise}
\leanok
\uses{thm:criticalPair}
Two raised arms never flat-join, a fork with an uncovered witness
joins, and no raise survives a twin-raise: the rung obstruction
that bars the difference from riding without a fork.
\end{theorem}

\begin{proof}
\leanok
Discharged in the Orig2 sources --- the card's axiom battery prints [propext, Quot.sound].
\end{proof}

\section{The cover and the class theorem}

\begin{theorem}[The cover lock]\label{thm:cover}
\lean{cover_top_of_fork, cover_pins_of_fork}
\leanok
\uses{def:fork}
At a fork the witness twin sits on a cover at the top of its run,
and the cover pins cascade: the seating census follows the cover
token.
\end{theorem}

\begin{proof}
\leanok
Discharged in the Orig2 sources --- the card's axiom battery prints [propext, Quot.sound].
\end{proof}

\begin{theorem}[The frozen cover rung]\label{thm:coverRung}
\lean{cover_rung_frozen}
\leanok
\uses{thm:cover}
The cover's rung is frozen along the class: no reversible move
can raise the covered card.
\end{theorem}

\begin{proof}
\leanok
Discharged in the Orig2 sources --- the card's axiom battery prints [propext, Quot.sound].
\end{proof}

\begin{theorem}[The fork-label class theorem]\label{thm:forkLabelClass}
\lean{fork_label_class}
\leanok
\uses{thm:cover, thm:rungObstruction}
Fork labels are class invariants: class-mate forks carry the same
label.  The two residues (cover token, one founding) are
discharged, so the class theorem is unconditional.
\end{theorem}

\begin{proof}
\leanok
Discharged in the Orig2 sources --- the card's axiom battery prints [propext, Quot.sound].
\end{proof}

\begin{theorem}[The pair founding row]\label{thm:twinPair}
\lean{twinPair_founds}
\leanok
\uses{thm:forkLabelClass}
The founding of the twin pair is a class invariant as well.
\end{theorem}

\begin{proof}
\leanok
Discharged in the Orig2 sources --- the card's axiom battery prints [propext, Quot.sound].
\end{proof}

\section{Fork-position rigidity}

\begin{definition}[Fork-position rigidity]\label{def:forkPosition}
\lean{ForkPositionRigid, ForkFoundRigid}
\leanok
\uses{def:position, def:fork}
All forks of one macro class share one Position (the
\emph{rigidity} claim) and one founding pattern (the \emph{found}
half).  Both halves are now unconditional: a fork class carries
at most one fork Position.
\end{definition}

\begin{theorem}[The pair half holds]\label{thm:forkPairClass}
\lean{ForkPairClass, forkPairClass, fork_under_cards}
\leanok
\uses{def:forkPosition}
The pair half of rigidity holds: pair suits found via the class
theorem, and at a non-king S-fork both twins sit on
opposite-family successor cards.
\end{theorem}

\begin{proof}
\leanok
Discharged in the Orig2 sources --- the card's axiom battery prints [propext, Quot.sound].
\end{proof}

\begin{theorem}[Rigidity reduces to the third suit]\label{thm:rigidReduce}
\lean{forkFoundRigid_iff_forkPositionRigid, forkPositionRigid_of_third, fork_one_position_of_rigid, position_eq_iff}
\leanok
\uses{def:forkPosition, thm:forkPairClass}
Found-rigidity and position-rigidity are equivalent; both follow
from the third-suit claim \texttt{ForkThirdRigid}, which is
therefore exactly what remained --- and is now closed.
\end{theorem}

\begin{proof}
\leanok
Discharged in the Orig2 sources --- the card's axiom battery prints [propext, Quot.sound].
\end{proof}

\section{The third-suit split}

\begin{theorem}[The third-suit reduction]\label{thm:thirdReduce}
\lean{ThirdRigid.forkThirdRigid_of_high}
\leanok
\uses{def:forkPosition}
The third-suit claim follows from its high-grade form
\texttt{ForkThirdRigidHigh}.
\end{theorem}

\begin{proof}
\leanok
Discharged in the Orig2 sources --- the card's axiom battery prints [propext, Quot.sound].
\end{proof}

\begin{theorem}[The non-king half]\label{thm:thirdNonKing}
\lean{ThirdHigh.forkThirdRigidHigh_of_not_king}
\leanok
\uses{thm:thirdReduce}
The high-grade third-suit rigidity holds when the fork witness is
not a king: the supported-card analysis confines the third suit.
\end{theorem}

\begin{proof}
\leanok
Discharged in the Orig2 sources --- the card's axiom battery prints [propext, Quot.sound].
\end{proof}

\begin{theorem}[The king slice]\label{thm:thirdKing}
\lean{ThirdHigh.ForkThirdRigidHighKing, ThirdHigh.forkThirdRigidHigh_of_kingSeated, ThirdHigh.forkThirdRigid_of_kingSeated}
\leanok
\uses{thm:kingExodus}
The king slice of the high-grade claim is reduced to the exodus
seating property: seated exodus kings close the gap.
\end{theorem}

\begin{proof}
\leanok
Discharged in the Orig2 sources --- the card's axiom battery prints [propext, Quot.sound].
\end{proof}

\begin{theorem}[The high third-suit split]\label{thm:thirdHigh}
\lean{ThirdRigid.ForkThirdRigidHigh, ThirdHigh.forkThirdRigidHigh_unconditional}
\leanok
\uses{thm:thirdReduce, thm:thirdNonKing, thm:thirdKing}
The high-grade third-suit rigidity, unconditional: the non-king
half plus the king slice (via the exodus seating) close the
split.
\end{theorem}

\begin{proof}
\leanok
\uses{thm:thirdNonKing, thm:thirdKing}
The non-king half confines the witness when it is not a king;
the king slice rides the exodus seating.
\end{proof}

\begin{theorem}[King exodus seating]\label{thm:kingExodus}
\lean{ThirdHigh.KingExodusSeated, ThirdHigh.kingExodusSeated_proved, ThirdHigh.KingKeep, ThirdHigh.king_keep_step, ThirdHigh.king_keep_transit}
\leanok
\uses{thm:cover, thm:criticalPair}
The exodus king is seated: a king founded at the deep end of a
class journey ends seated in the tableau, never stranded.  The
proof rides the \texttt{KingKeep} invariant --- a king is founded
or bottoms a hidden-free pile --- which is closed under every
reversible step and survives the whole witness play.
\end{theorem}

\begin{proof}
\leanok
\uses{thm:coverRung, thm:forkPairClass}
The king trace: a king's \texttt{foundToTab} landing base is an
empty anchor, the trace is preserved by every reversible letter,
and transitivity carries it along the witness play.  The last
open row of the rigidity arc, closed.
\end{proof}

\chapter{The fiber hierarchy and the arm split}
\label{chap:fiber}

\section{The dichotomy family}

\begin{definition}[The dichotomy family]\label{def:dichotomy}
\lean{ArmDichotomyFor, ArmDichotomy, ArmDichotomySW, ArmDichotomyRung, SatEqualFork, SatEqualCargoSW}
\leanok
\uses{def:fibers, def:ssaturated}
For a relation $R$ on states, \texttt{ArmDichotomyFor} $R$ says:
saturated same-orbit mates are either $R$-related or a fork lives
in their class.  The cargo, flat (SW) and rung readings give the
three arms; the equal-rung readings give the fork-equality
family.
\end{definition}

\begin{theorem}[The arm split]\label{thm:armSplit}
\lean{armDichotomy_iff_split, satEqualFork_of_armDichotomy}
\leanok
\uses{def:dichotomy, thm:satDifferFork}
The generic arm dichotomy splits exactly into the differing-rung target
and the equal-rung half.  The cargo-fiber form of the equal-rung half is
refuted by the draw boundary, so only the generic split and its
\texttt{SatDifferFork} projection remain; the SW flat-fiber repair is the
live replacement.
\end{theorem}

\begin{proof}
\leanok
Discharged in the Orig2 sources --- the card's axiom battery prints [propext, Quot.sound].
\end{proof}

\begin{theorem}[The flat-fiber dichotomy]\label{thm:swDichotomy}
\lean{armDichotomySW_iff_split, armDichotomySW_to_satDifferFork, forkless_found_unique_SW}
\leanok
\uses{def:dichotomy}
The flat-fiber (SW) dichotomy splits the same way, implies the
differing-runk target, and supplies forkless foundation
uniqueness in the flat fiber.
\end{theorem}

\begin{proof}
\leanok
Discharged in the Orig2 sources --- the card's axiom battery prints [propext, Quot.sound].
\end{proof}

\section{The class-global invariant}

\begin{definition}[The flat-fiber vertex]\label{def:satFlatFiberFork}
\lean{SatFlatFiberFork, SatCargoFiberFork, ForklessFlatNormalWF}
\leanok
\uses{def:dichotomy}
The class-global invariant: two saturated class-mates in
different flat fibers force a fork somewhere in the class.  The
draw-free cargo variant is recorded only to fence the wrong
vertex --- it is false in a live draw zone.  The forkless form
says all saturated class-mates of a forkless class lie in one
flat fiber.
\end{definition}

\begin{theorem}[The flat vertex discharges the target]\label{thm:flatVertex}
\lean{not_bareFlat_of_rung_ne, satDifferForkWF_of_satFlatFiberFork, satFlatFiberFork_of_armDichotomySW, satDifferForkWF_of_armDichotomySW, rung_diff_forces_fork, forklessFlatNormalWF_of_armDichotomySW}
\leanok
\uses{def:satFlatFiberFork, thm:swDichotomy}
A rung difference rules out the flat fiber, so the invariant
discharges the WF-gated target; and the SW arm dichotomy supplies
the invariant directly --- its non-flat branch is exactly the
invariant.  The campaign's vertex is the flat fiber, not the
draw-fragile cargo fiber.
\end{theorem}

\begin{proof}
\leanok
Discharged in the Orig2 sources --- the card's axiom battery prints [propext, Quot.sound].
\end{proof}

\begin{theorem}[The ceiling bridge]\label{thm:ceiling}
\lean{forkless_rungs_agree, forkless_found_unique, forkless_found_unique_wf}
\leanok
\uses{thm:satDifferFork}
The campaign payoff, conditional on the target: forkless
saturated classes have rung-agreeing endpoints, and foundations
are unique in a forkless class.  One \texttt{elim} from the
target, and the rung-to-found congruence finishes.
\end{theorem}

\begin{proof}
\leanok
Discharged in the Orig2 sources --- the card's axiom battery prints [propext, Quot.sound].
\end{proof}

\section{The census rows}

\begin{theorem}[The rung census]\label{thm:rungCensus}
\lean{shuffle_rung_census, found_exodus, rung_diff_exhibit}
\leanok
\uses{def:orbit}
The rung ledger: class journeys census their rung steps, the
differing card is pulled in LIFO order (the exodus), and
differing rungs are exhibited per suit.
\end{theorem}

\begin{proof}
\leanok
Discharged in the Orig2 sources --- the card's axiom battery prints [propext, Quot.sound].
\end{proof}

\chapter{The boundary cases}
\label{chap:boundary}

\section{The equal-rung boundary}

\begin{theorem}[The equal-rung boundary]\label{thm:vcBoundary}
\lean{vc_equalRung_boundary, vc_armDichotomy_fork}
\leanok
\uses{def:dichotomy}
Two saturated class-mates with the same rung vector
$(3,0,4,3)$, joined by a twelve-leg reversible witness journey,
and not cargo-equivalent.  The cargo dichotomy's first branch is
refuted here, so the residue must supply a class-wide fork.
\end{theorem}

\begin{proof}
\leanok
Discharged in the Orig2 sources --- the card's axiom battery prints [propext, Quot.sound].
\end{proof}

\section{The draw boundary}

\begin{theorem}[The draw boundary]\label{thm:drawBoundary}
\lean{satEqualFork_draw_false, armDichotomy_draw_false}
\leanok
\uses{def:dichotomy}
The all-non-king board (foundations full below the kings, empty
tableau, the four kings split between stock and waste) is
draw-only, and the four-draw rotation makes two saturated
same-orbit mates in different cargo fibers.  Hence the equal-runk
fork claim and the cargo dichotomy are both refuted outright.
\end{theorem}

\begin{proof}
\leanok
Discharged in the Orig2 sources --- the card's axiom battery prints [propext, Quot.sound].
\end{proof}

\begin{theorem}[The amended dichotomy]\label{thm:amended}
\lean{armDichotomyRung_iff_satDifferFork}
\leanok
\uses{def:dichotomy, thm:satDifferFork}
The amended dichotomy --- first branch the rung ledger equality
instead of the cargo fiber --- is exactly the differing-runk
target: two statements, one case split apart.
\end{theorem}

\begin{proof}
\leanok
Discharged in the Orig2 sources --- the card's axiom battery prints [propext, Quot.sound].
\end{proof}

\section{The journey-local dip}

\begin{theorem}[The journey-local dip is false]\label{thm:dipFalse}
\lean{ForkFree.journey_local_dip_false, ForkFree.not_satDip, ForkFree.not_satJourneyFork}
\leanok
\uses{def:fork}
Two Plan-4 saturated boards with rungs $(5,0,4,2)$ and
$(5,0,2,5)$ are joined by a fifteen-state, fourteen-move witness
play that avoids all eight fork states of the class.  The
journey-local residues are refuted; the surviving residue is the
class-global target.
\end{theorem}

\begin{proof}
\leanok
Discharged in the Orig2 sources --- the card's axiom battery prints [propext, Quot.sound].
\end{proof}

\begin{theorem}[The minimal path crosses the fork]\label{thm:minPath}
\lean{SatDifferMin9P.minimal_play_crosses_fork_plan4}
\leanok
\uses{def:fork, thm:dipFalse}
The same Plan-4 boards admit a twelve-move, nine-bracket witness
play whose sixth state is a $(3,0,2,2)$ fork --- the first fork
on the path --- while the fourteen-move fork-free play costs
eleven brackets.  The fork-free route is not bracket-minimal: it
pays one extra pull-and-raise pair.  The kernel-checked
representative of the finite-model census; the universal lower
bound (every nine-bracket path crosses a fork) remains in the
census.
\end{theorem}

\begin{proof}
\leanok
Discharged in the Orig2 sources --- the card's axiom battery prints [propext, Quot.sound].
\end{proof}

\begin{theorem}[The first-passage ledger]\label{thm:firstPassage}
\lean{SatDifferCensus.step_rung_classification, SatDifferCensus.firstPassage_exists, SatDifferCensus.firstRise_or_firstFall_of_rung_ne}
\leanok
\uses{def:orbit}
Every suit-crossing of a witness journey has a first passage, and
each decomposes into a first rise or a first fall: the ledger the
campaign's route reads.
\end{theorem}

\begin{proof}
\leanok
Discharged in the Orig2 sources --- the card's axiom battery prints [propext, Quot.sound].
\end{proof}

\chapter{The campaign targets}
\label{chap:campaign}

\section{The main target}

\begin{theorem}[Differing rungs force a fork]\label{thm:satDifferFork}
\lean{SatDifferFork, SatDifferForkWF}
\leanok
\uses{def:fork, def:ssaturated}
Two S-saturated class-mates whose rung vectors disagree force a
fork somewhere in the class: formally, for saturated same-orbit
$u, v$ with $\mathrm{rungAt}(u, \sigma) \ne \mathrm{rungAt}(v,
\sigma)$ for some suit $\sigma$, there is a class-mate $F$ of
$u$ with \texttt{ForkWitness} $F\,x$.  The contrapositive reads:
forkless classes have rung-agreeing saturated endpoints.
\end{theorem}

\begin{proof}
\uses{thm:flatVertex, thm:bridge, thm:forkLabelClass, thm:rungCensus, thm:firstPassage, thm:dipFalse, thm:minPath}
Open --- no formal proof yet, on two fronts.  The vertex front:
the flat-fiber invariant discharges the WF-gated form, and the SW
arm dichotomy supplies the invariant, so the gap is the SW
dichotomy itself.  The valley front: the drain is complete (fork
existence from the asymmetry is unconditional) and the bridge
from the bottom asymmetry to the target is proved from the
corrected residual; the campaign hangs from two named finite
residuals --- valley elimination at the minimality gate (the
bracket existence, after the normalizer and the read-invariance
induction) and the split block (the post-cover analysis).  The saturation grade is
load-bearing at the endpoints: it bars the cover's own raise and
the covered card's cascade license.  The journey-local versions
are refuted, the minimal-path boundary certifies the detour's
cost, and the payoff is the ceiling bridge.
\end{proof}

\section{The valley route}

The live attack front, now some twenty-five cards strong, on two
converging routes.  The crossing route runs through the normalizer
(the commutation kit, the read-invariance induction, the bracket
existence at the minimality gate); the bridge route runs through
the corrected split residual (the live pair plus the block).  The
refutations on the way --- the mountain form, the ungated
locator, the transport rows, the binder-scope split residual ---
were each corrected into a sharper statement.

\begin{theorem}[Valley-at-a-fork]\label{thm:valleyForm}
\uses{def:fork}
The surviving normal form: every unit-shortest, lexicographic or
bracket-minimal witness play between saturated states is
componentwise a valley --- every turn state is a fork, and the
last pull and first raise at the bottom are twins.  The mountain
form is refuted in general (the raise card is buried; only a
twin-suit pull opens the slide seat) but true below the fork,
fork-to-fork, and on genuinely forkless classes.  Probe-verified
by the finite-model census; the Lean statements are named but
not yet written.
\end{theorem}

\begin{theorem}[The valley-turn card]\label{thm:valleyTurn}
\lean{ValleyStep.valley_cancel_same, ValleyStep.valley_cancel_at_fork_twin, ValleyStep.fork_valley_cancel_of_pull, ValleyStep.fork_first_raise_eq, ValleyStep.fork_last_pull_eq, ValleyNoPull.nextUp_of_pull_of_ne}
\leanok
\uses{def:fork}
The same-card valley turn cancels, and at a fork the turn
cancels at the twin; the first raise and the last pull at a fork
are pinned to the witness; and a pull never re-opens the next-up
flag.  The promoted valley-turn rows, in the library.
\end{theorem}

\begin{proof}
\leanok
Discharged in the Orig2 cards (green, axiom battery [propext, Quot.sound]).
\end{proof}

\begin{theorem}[The valley-elimination kit]\label{thm:valleyKit}
\lean{ValleyElim.crossing_cancel, ValleyElim.crossing_block_delete, ValleyElim.valley_elim_pull_raise_same, ValleyElim.valley_exists}
\leanok
\uses{def:orbit}
The four-letter crossing block (pull $d$, pull $c$, raise $d$,
raise $c$) returns to its origin, so a contiguous crossing block
deletes and the play gets strictly shorter; same-card turns
eliminate unconditionally; and every witness play between
saturated endpoints with differing rungs has a valley split, free
letters allowed at the turn.
\end{theorem}

\begin{proof}
\leanok
Discharged in the Temp kit (green, axiom battery [propext, Quot.sound]).
\end{proof}

\begin{theorem}[Valley elimination, no forks]\label{thm:valleyTarget}
\lean{ValleyElim.ValleyElimNonfork, ValleyElim.ShortestPlayExists, ExistKit.ValleyElimNonforkMin}
\leanok
\uses{thm:valleyKit, thm:valleyTurn}
The lead target: crossing valleys can be eliminated without
ever introducing a fork, and a shortest play exists between any
two states.  The ungated crossing locator was refuted and
corrected (the partner locator); the elimination is reduced to
the bracket deletion at the minimality gate.
\end{theorem}

\begin{proof}
\uses{thm:normalizer, thm:minimalityGate}
The commuted brackets delete, the read-invariance induction
carries the letters past the normalizer's moves, and the
remaining obligation is the bracket existence at a
minimum-length play.  Open --- no formal proof yet.
\end{proof}

\begin{theorem}[The crossing normalizer]\label{thm:normalizer}
\lean{NonforkKit.RaiseSwaps, NonforkKit.bracket_play_of_swaps, NonforkKit.bracket_delete_gen, Bubble.DiamondAt, Bubble.Commutes, LocatorCorrected.ValleyCrossingPartnerLocator, LocatorCorrected.CrossingBracketNormalizer, LocatorCorrected.normalizer_shortens, Elim.bracket_delete_of_commuted, Elim.bubble_bracket_shortens}
\leanok
\uses{thm:valleyKit}
The commutation kit, promoted: raise sequences commute past
unrelated letters (the slide rows are now free in the
normalizer), commuted brackets delete, and the bubble shortens.
The crossing locator was refuted and corrected --- the partner
locator finds the matching pull and raise, and the corrected
normalizer shortens the play.  The ungated transport rows were
refuted: the commuting normalizer needs the read-invariance
side conditions.
\end{theorem}

\begin{proof}
\leanok
Discharged in the Orig2 cards (green, axiom battery [propext, Quot.sound]).
\end{proof}

\begin{theorem}[The read-invariance induction]\label{thm:readInvariance}
\lean{R12.canRaise_preserve_of_pile, R12.canRaise_survives_raise, R12.canRaise_survives_pull, R12.canRaise_survives_slide}
\leanok
\uses{thm:decidable}
CanRaise survives a raise in another suit, a pull, and a
slide: the read of a card's raisability is stable under exactly
the letters the normalizer moves.  With this the commutation
half is complete; only the bracket existence remains.
\end{theorem}

\begin{proof}
\leanok
Discharged in the Orig2 card (green, axiom battery [propext, Quot.sound]).
\end{proof}

\begin{theorem}[The minimality gate]\label{thm:minimalityGate}
\lean{ExistKit.ReadBracketElimMin, ExistKit.ValleyForkAtMin, ExistKit.readBracketElimMin_of_readBracket, ExistKit.bracket_of_samecard_window}
\leanok
\uses{thm:readInvariance, thm:normalizer}
The ungated bracket residuals are refuted --- a four-letter
loop play (pull, slide, slide back, raise) is a cycle whose
only bracket fails the read condition, exactly the kind of play
that cannot occur at minimum length.  The corrected residuals
carry the minimality gate: the valley fork occurs already at a
minimum-length play.
\end{theorem}

\begin{proof}
\leanok
Discharged in the Orig2 cards (green, axiom battery [propext, Quot.sound]).
\end{proof}

\begin{theorem}[The reduction]\label{thm:valleyReduction}
\lean{ValleyElim.satDifferForkWF_of_valleyElimNonfork}
\leanok
\uses{thm:valleyTarget, def:fork, def:ssaturated}
Valley elimination without forks, plus a shortest play,
discharges the WF-gated campaign target.  Together with the drain
skeleton, the target now hangs from named scratch-pen inputs
rather than an unpriced route.
\end{theorem}

\begin{proof}
\leanok
Discharged in the Temp kit (green, axiom battery [propext, Quot.sound]).
\end{proof}

\begin{theorem}[The drain-to-fork skeleton]\label{thm:drainSkeleton}
\lean{Drain.drain_to_fork, Drain.drain_terminal_fork, Drain.drain_terminal_forkState, Drain.raise_step_preserves, Drain.rungSum_le_52, Drain.drainTerminal_decidable, ValleyDrain.canRaise_of_raise_other, ValleyDrain.not_canRaise_of_raise_other, ValleyDrain.canRaise_of_raise_high}
\leanok
\uses{thm:safePath, def:fork}
The skeleton of the descent, promoted into the library: raise
the live cards while the blocked twin stays blocked (an
unrelated or higher-ranked raise cannot be the cover),
terminate by rank monotonicity (the rung sum is bounded by 52),
and the terminal state is a fork outright.  Every named input is
now discharged --- the transport by the safe transport, the
regulation by the one-slide live card, and the step by the
closed form --- completing the drain below.
\end{theorem}

\begin{proof}
\leanok
Discharged in the Orig2 card (green, axiom battery [propext, Quot.sound]).
\end{proof}

\begin{theorem}[The safe cargo path and the transport]\label{thm:safePath}
\lean{SafePathFiber.safe_cargo_path_unconditional, SafePathFiber.twin_top_of_cover, SafePathFiber.cover_slide_exists, SafeTransport.canRaiseS_of_canRaiseS_raise_high, SafeTransport.canRaiseS_twin_of_raise_high}
\leanok
\uses{thm:decidable, def:fork}
From twin S-liveness and the vanilla block, the twin is the pile
top and one slide --- the cover card onto the twin's seat --- is
legal, reversible, and safe: the safe cargo path is unconditional
and a single slide.  Raising a card of rank at least the twin's,
in a different suit and not the twin itself, preserves both the
slide and the twin's S-liveness; the side condition is necessary
(a full-deck counterexample card certifies it), and the drain
never raises the twins.  This discharges the skeleton's first
named input.
\end{theorem}

\begin{proof}
\leanok
Discharged in the Orig2 cards (green, axiom battery [propext, Quot.sound]).
\end{proof}

\begin{theorem}[The one-slide live regulation]\label{thm:oneSlideLive}
\lean{OneSlide.one_slide_of_blocked_live, OneSlide.twin_top_of_cover, OneSlide.cover_slide_exists, OneSlide.cover_unique}
\leanok
\uses{thm:decidable, def:fork}
The live-card regulation: the twin tops the cover pile, the
cover is unique, and no twin sits below the blocked card --- the
one-slide frame confines every live card at the drain.  This
discharged the skeleton's regulation input.
\end{theorem}

\begin{proof}
\leanok
Discharged in the Orig2 card (green, axiom battery [propext, Quot.sound]).
\end{proof}

\begin{theorem}[The drain, complete]\label{thm:drainComplete}
\lean{DrainClose.fork_exists_of_asymmetry, DrainClose.drain_to_fork_discharged, DrainClose.slide_frame, DrainClose.floor_slide, DrainClose.floor_raise, DrainClose.tabToFoundR_of_canRaise}
\leanok
\uses{thm:drainSkeleton, thm:safePath, thm:oneSlideLive}
Fork existence is unconditional: from a live card whose twin is
S-live but vanilla-blocked, with the rank floor, a fork exists
in the class.  The slide frame gives the four exact pile rows of
the cover slide; the rank floor holds on both halves (the slide
by cargo-reachability, the raise by the flat characterization);
the side-residuals vanished --- every vanilla raise is
reversible, and blocked cards are handled by slide-and-raise.
Termination is the measure $52 - \mathrm{rungSum}$; the terminal
branch is the fork.
\end{theorem}

\begin{proof}
\leanok
Discharged in the Orig2 card (green, axiom battery [propext, Quot.sound]).
\end{proof}

\begin{theorem}[The corrected split residual]\label{thm:splitResidual}
\lean{SplitKit.SplitResidualData, SplitKit.splitResidualData_iff_parts, SplitKit.splitResidualData_of_livePair_block, SplitKit.splitTwin_of_livePair_block, SplitKit.splitFloor_of_livePair_twin, SplitKit.bottomDip_of_splitResidualData, SplitKit.not_splitResidual_of_split}
\leanok
\uses{thm:valleyKit, thm:valleyTurn}
The corrected residual.  The original statement was false as
literally written --- its split binders never occurred in the
hypotheses, so it demanded the twin equation for arbitrary
cards, refuted by machine --- and the correction supplies the
split equations as hypotheses.  The corrected form decomposes
three ways (the twin equation, the block, the rank floor), and
the two-way reduction proves the twin equation and the floor
outright from the live pair plus the block.  The corrected dip
row feeds the bridge.
\end{theorem}

\begin{proof}
\uses{thm:splitMechanism}
The residual is the live pair plus the block; the mechanism
chain below closes both halves in pieces.  Open --- no formal
proof yet.
\end{proof}

\begin{theorem}[The split mechanism]\label{thm:splitMechanism}
\lean{LiveKit.splitLivePair_of_valleyElim, LiveKit.splitLivePair_of_forkAtRaise, Landing.landingLivePair_of_twinLivePair_F, Landing.splitLivePair_of_landing_parts, Landing.ThirdSuitBracket, Landing.thirdSuitBracket_iff_landingThirdSuitFree, BlockKit.PostCover, BlockKit.splitBlock_iff_postCover, BlockKit.NoEmptyTail, BlockKit.TailCovers, Mech.NoMover, Mech.mover_or_last_mover}
\leanok
\uses{thm:valleyKit}
The mechanism chain: the live pair follows from valley
elimination outright (and from a fork at the raise); the landing
live pair equals the split live pair; the block is a
post-cover --- the free tail covering --- reduced to
no-empty-tail plus tail covers, with the pile half closed; and
the mover analysis pins the seat changes.  The last step is
the third-suit bracket.
\end{theorem}

\begin{proof}
\leanok
Discharged in the Orig2 cards (green, axiom battery [propext, Quot.sound]).
\end{proof}

\begin{theorem}[The asymmetry-production bridge]\label{thm:bridge}
\lean{Bridge.satDifferForkWF_of_bottom_asymmetry, Bridge.forkExistence_of_drainInputs, Bridge.bottom_asymmetry_of_split_data, Bridge.bottomAsymmetry_of_bottomDip, Bridge.valley_bottom_no_twin_block, Bridge.valley_bottom_twins_both_raisable, Bridge.no_forkWitness_at_pull_landing, Bridge.blocked_pull_needs_free_tail}
\leanok
\uses{thm:splitResidual, thm:drainComplete}
The bridge from the residual to the target: the split residual
implies the bottom dip, the dip implies the bottom asymmetry,
and with the drain complete the asymmetry discharges the
WF-gated target.  The no-go rows confine the residual's shape:
the raise card is reversible hence never the blocked twin, at
the pull's landing both twins are raisable so the cover must
come from a free slide after the pull, and the pull-landing
bottom is never a fork.
\end{theorem}

\begin{proof}
\leanok
Discharged in the Orig2 card (green, axiom battery [propext, Quot.sound]).
\end{proof}

\section{The rigidity arc, closed}

\begin{theorem}[Fork-position rigidity]\label{thm:forkPositionTarget}
\lean{ForkPositionRigid, ThirdHigh.forkPositionRigid_unconditional, ForkFoundRigid, ThirdHigh.forkFoundRigid_unconditional}
\leanok
\uses{thm:rigidReduce, thm:thirdHigh}
All forks of a class share one Position, and the found half holds
with it: a fork class carries at most one fork Position.
Unconditional.
\end{theorem}

\begin{proof}
\leanok
\uses{thm:rigidReduce, thm:thirdHigh}
By the reduction, rigidity follows from the high third-suit
claim, now closed unconditionally: the non-king half plus the
king slice via the exodus seating.
\end{proof}
